\section{Experiments}
\label{sec:experiments}

\subsection{Setup}
\paragraph{Datasets.}
\textbf{RepCount-A}~\cite{hu2022transrac} contains 1{,}041 videos of fitness
and daily activities with per-cycle start/end annotations, split into
758/131/152 train/validation/test videos. Median cycle lengths range from
0.85\,s to 4.75\,s (5th--95th percentile of per-video medians), and 20\% of the
videos vary their cycle length by more than 25\%.
\textbf{UCFRep}~\cite{zhang2020context} annotates 526 UCF101 videos with cycle
boundaries (421/105 split); following Context and ESCounts we count inside the
annotated repetitive segment. Because UCF101 clips of one group are cut from the
same source video, model selection uses a development split of the 421 training
videos made of whole groups (331/90). We do not retrain on the 90 held-out
videos, so \method uses fewer UCFRep training videos than prior work.
\textbf{Countix}~\cite{dwibedi2020repnet} consists of Kinetics clips with counts.
We use the train/val/test split of ESCounts and retrieve 98\% of its videos
(4{,}224/1{,}374/2{,}403); exclusions (unresolvable or missing clips, exact
duplicate rows) are listed in the supplement. Countix has no cycle boundaries,
so density targets place the counted cycles uniformly over the repetition span,
as in ESCounts.
\textbf{QUVA}~\cite{runia2018realworld} (100 videos) is used only for zero-shot
evaluation of RepCount-A models. We verified every count and boundary against an
independent copy of the official annotations (zero mismatches on RepCount-A and
UCFRep).

\paragraph{Metrics and protocol.}
We report the normalised error
$\mathrm{MAE}=\frac{1}{N}\sum_i |\mathrm{round}(\hat c_i)-c_i|/(c_i+0.1)$ and
the off-by-one accuracy
$\mathrm{OBO}=\frac{1}{N}\sum_i\mathbb{1}[|\mathrm{round}(\hat c_i)-c_i|\le1]$
over videos with a positive count, exactly as in the released ESCounts
evaluation code; with this protocol the released ESCounts checkpoint obtains
0.215/0.556 on RepCount-A test in our pipeline, against its published
0.213/0.563. One training run on the union of the three training sets yields one
model per benchmark, selected on that benchmark's validation split (UCFRep: on
the development split) and evaluated once on its test split; no test video is
used for training or selection. We train five seeds (ablations: three) and
report mean$\pm$standard deviation; differences are tested with a paired
video-level bootstrap.

\paragraph{Implementation.}
The video encoder is the frozen ST-MAE ViT-B/16~\cite{feichtenhofer2022masked}
(Kinetics-400 self-supervised; the encoder of ESCounts), with shorter side 224, a
centre crop and ImageNet normalisation; tokens are taken after blocks 8 and 12.
The MLLM is Qwen3-VL-8B-Instruct~\cite{qwen3vl} in bf16, frozen, run once per
video with greedy decoding of at most 48 tokens. The decoder has width 256 and
6.1M trainable parameters; we use AdamW (learning rate $2\times10^{-4}$, weight
decay 0.05, cosine schedule with 3 warm-up epochs), batches of 8 whole videos,
50 epochs over the balanced union, random temporal crops covering at least 60\%
of a video with probability 0.5, pyramid shift with probability 0.5, and bf16
mixed precision. All experiments run on two RTX 4090 GPUs; caching the MLLM
outputs for all 9.8k videos takes about 2.7 hours, and one training run about
2 hours.

\subsection{Comparison with the state of the art}
\Cref{tab:main} \TODO{results discussion}.

\subsection{Ablations}
\Cref{tab:ablation} \TODO{ablation discussion: pyramid, overlap averaging,
objective, depth}.

\subsection{Generalisation and robustness}
\TODO{QUVA zero-shot, RepCount$\to$UCFRep/Countix, tempo-warp stress test.}
